%& -shell-escape
% !TeX root = thesis-abstract.tex
% !TeX encoding = UTF-8
% !TeX program = XeLaTeX

\documentclass{.local/lib/classes/ndvr-super-article-standalone}

\begin{document}

    Видео содержит в себе несколько видов данных. 
    Выделяют четыре вида:\footnote{
        Работа Смитона Techniques used and open challenges 
        to the analysis, indexing and retrieval of digital video. 
            Information Systems, 32(4):545–559, 2006.
    }
    \begin{itemize}
        \item метаданные — заголовок, автор и описание;
        \item звуковая дорожка;
        \item тексты, полученные при помощи технологии 
            оптического распознавания символов (OCR) 
            или при помощи технологий распознавания речи;
        \item визуальная информация кадров видео. 
    \end{itemize}

    Источники данных разной природы 
    не являются независимыми. Они тесно связаны и дополняют 
    друг друга, позволяя воссоздать полную 
    картину происходящего.
    Связь различных видов данных представлена
    на рисунке \ref{fig:semantic-uml}.

    \begin{figure}[H]
        \includegraphics[width=\textwidth]
            {.local/vec/out/pdf/video-model-uml-semantic.pdf}
        \caption{%
            Логическая модель видео,\\
            выраженная диаграммой классов UML 2.0%
        }
        \label{fig:semantic-uml}
    \end{figure}

    Часто видео рассматривается 
    как~последовательность кадров.
    Это связано с особенностями физического представления видео.
    Модель физического представления видео показана  
    на рисунке \ref{fig:phisical-uml}.

    \begin{figure}[p]
        \includegraphics[width=\textwidth]
            {.local/vec/out/pdf/video-model-uml-logical.pdf}
        \caption{%
            Логическая модель физического представления видео,\\
            выраженная диаграммой классов UML 2.0%
        }
        \label{fig:phisical-uml}
    \end{figure}
    
\end{document}
